\documentclass[11pt]{article}
\usepackage[a4paper,margin=1in]{geometry}
\usepackage{fontspec}
\usepackage[boldfont=false]{xeCJK}
\setCJKmainfont{FandolSong-Regular.otf}
\usepackage{amsmath}
\usepackage{amssymb}
\usepackage{array}
\usepackage{booktabs}
\usepackage{graphicx}
\usepackage{adjustbox}
\usepackage{enumitem}
\setlist[itemize]{topsep=2pt,partopsep=0pt,parsep=0pt,itemsep=2pt,leftmargin=1.4em}
\setlist[enumerate]{topsep=2pt,partopsep=0pt,parsep=0pt,itemsep=2pt,leftmargin=1.6em}
\usepackage{parskip}
\usepackage{fvextra}
\fvset{breaklines=true,breakanywhere=true,fontsize=\small}
\setlength{\parindent}{0pt}

\begin{document}

\begin{center}
  {\LARGE\bfseries Programming}\\[0.35em]
  {\large A Level Computer Science}
\end{center}
\vspace{0.6em}

\section*{Programming basics}

\subsection*{From design to code}

You should be able to turn a design — a \textbf{flowchart} 流程图 or \textbf{structured English} 结构化英语 — into \textbf{pseudocode} 伪代码, and then into a real language:

\begin{enumerate}
\item find the \textbf{variables} 变量 and their \textbf{data types} 数据类型.

\item turn input/output boxes into \texttt{INPUT} / \texttt{OUTPUT}.

\item turn decision diamonds into \texttt{IF...ELSE...ENDIF} (or \texttt{CASE}).

\item turn loop arrows into \texttt{WHILE}, \texttt{REPEAT...UNTIL}, or \texttt{FOR}.

\item turn process boxes into assignments or calculations.

\item check by tracing a small input.

\end{enumerate}

\subsection*{Constants and variables}

A \textbf{constant} 常量 holds a value that \textbf{never changes}; a variable holds one that may change. Declare them with a type:

\begin{Verbatim}
CONSTANT Pi ← 3.14159
DECLARE Radius : REAL
DECLARE Area : REAL

Radius ← 5
Area ← Pi * Radius * Radius
\end{Verbatim}

Use constants for fixed values that recur (\texttt{Pi}, \texttt{MaxScore}); they make code clearer and easy to change in one place.

\subsection*{Assignment and expressions}

Use \texttt{←} for \textbf{assignment} 赋值:

\begin{Verbatim}
Total ← Total + 1
Average ← Sum / Count
\end{Verbatim}

Expressions use \textbf{operators} 运算符:

\begin{itemize}
\item arithmetic \texttt{+ - * /}, plus \texttt{DIV} (integer division) and \texttt{MOD} (remainder): \texttt{7 DIV 2 = 3}; \texttt{7 MOD 2 = 1}.

\item comparisons \texttt{=}, \texttt{\textless{}\textgreater{}}, \texttt{\textless{}}, \texttt{\textgreater{}}, \texttt{\textless{}=}, \texttt{\textgreater{}=}.

\item logic \texttt{AND}, \texttt{OR}, \texttt{NOT}.

\end{itemize}

\textbf{Precedence} 优先级 (highest to lowest): \texttt{NOT} → \texttt{* / DIV MOD} → \texttt{+ -} → comparisons → \texttt{AND} → \texttt{OR}. Use brackets when unsure.

\subsection*{Input and output}

\begin{Verbatim}
OUTPUT "Enter your name:"
INPUT Name
OUTPUT "Hello, ", Name
\end{Verbatim}

\subsection*{Built-in functions and library routines}

Many tasks have ready-made \textbf{library routines} 库例程, so you need not write them:

\begin{itemize}
\item string: \texttt{LENGTH(s)}, \texttt{LEFT(s, n)}, \texttt{RIGHT(s, n)}, \texttt{MID(s, start, len)}, \texttt{UCASE(s)}, \texttt{LCASE(s)}.

\item numeric: \texttt{INT(x)}, \texttt{ROUND(x)}, \texttt{ABS(x)}, \texttt{MOD(a, b)}, \texttt{RANDOM()}.

\item conversion: \texttt{STR(x)} (number → string), \texttt{VAL(s)} (string → number).

\end{itemize}

Use the exact names from the question paper's reference list.

\section*{Selection}

\textbf{Selection} 选择 chooses which steps run.

\begin{Verbatim}
IF age >= 18 THEN
    OUTPUT "Adult"
ELSE
    OUTPUT "Minor"
ENDIF
\end{Verbatim}

For more than two cases you can use a \textbf{nested} 嵌套 IF, but deep nesting is hard to read — a \texttt{CASE} is cleaner when testing one value against several options:

\begin{Verbatim}
CASE OF Grade
    "A": OUTPUT "Excellent"
    "B": OUTPUT "Good"
    OTHERWISE: OUTPUT "Try again"
ENDCASE
\end{Verbatim}

Cambridge \texttt{CASE} allows single values, value lists (\texttt{1, 2, 3:}), and ranges (\texttt{1 TO 5:}).

\par\medskip
\noindent\begin{minipage}{\linewidth}
\centering
\adjustbox{max width=\linewidth, max totalheight=0.42\textheight}{%
  \includegraphics{/Users/joschka/Desktop/Programming/eXercise/output/eXam/handouts/a_level_computer_science/assets/11-p288-case-flowchart.pdf}}\\[0.3em]
\emph{A CASE statement runs the branch that matches the value}
\end{minipage}\par\medskip

\section*{Iteration}

\textbf{Iteration} 迭代 repeats a block. Three loops differ in how many times the body runs.

\subsection*{Count-controlled (FOR) loop}

A \textbf{count-controlled loop} 计数循环 — use it when you know how many times to repeat:

\begin{Verbatim}
FOR i ← 1 TO 10
    OUTPUT i
NEXT i
\end{Verbatim}

A \texttt{STEP} can change the count (e.g. \texttt{FOR i ← 10 TO 1 STEP -1}). Best for a fixed number of repeats or processing each element of an \textbf{array} 数组.

\subsection*{Pre-condition (WHILE) loop}

A \textbf{pre-condition loop} 前测循环 tests the condition \textbf{before} each pass, so it may run \textbf{zero} times:

\begin{Verbatim}
WHILE total < 100 DO
    INPUT n
    total ← total + n
ENDWHILE
\end{Verbatim}

\subsection*{Post-condition (REPEAT...UNTIL) loop}

A \textbf{post-condition loop} 后测循环 tests the condition \textbf{after} each pass, so it always runs \textbf{at least once}:

\begin{Verbatim}
REPEAT
    INPUT password
UNTIL password = correctPassword
\end{Verbatim}

\subsection*{Choosing the right loop}

\begin{itemize}
\item count known up front → \textbf{FOR}.

\item may need zero passes → \textbf{WHILE}.

\item always at least one pass → \textbf{REPEAT...UNTIL}.

\end{itemize}

Justify your choice by whether the count is known and whether the body must run at least once. A typical question gives a scenario ("ask for a password until correct, but always ask at least once") and asks which loop fits.

\section*{Procedures and functions}

\textbf{Structured programming} 结构化编程 builds a program from small named \textbf{subroutines} 子程序, each with one job.

\subsection*{Procedure}

A \textbf{procedure} 过程 is a named block that \textbf{does an action}; it may take \textbf{parameters} 参数 but does \textbf{not return a value}.

\begin{Verbatim}
PROCEDURE Greet(name : STRING)
    OUTPUT "Hello, ", name
ENDPROCEDURE

CALL Greet("Ada")
\end{Verbatim}

\subsection*{Function}

A \textbf{function} 函数 is like a procedure but it \textbf{returns a value} that becomes part of an expression.

\begin{Verbatim}
FUNCTION Square(x : INTEGER) RETURNS INTEGER
    RETURN x * x
ENDFUNCTION

result ← Square(5) + 1     // result = 26
\end{Verbatim}

Use a \textbf{procedure} when the subroutine performs an action; use a \textbf{function} when it computes a value for the caller.

\subsection*{Parameters}

A \textbf{parameter} is a variable a subroutine declares to receive input; the values the caller supplies are \textbf{arguments} 实参. Two ways to pass them:

\begin{itemize}
\item \textbf{pass by value} 传值 — the routine gets a \textbf{copy}; changes inside it do not affect the caller. Use for inputs it only reads.

\item \textbf{pass by reference} 传引用 — the routine gets a \textbf{reference} to the caller's variable; changes \textbf{do} affect the caller. Use when it must update a parameter.

\end{itemize}

\begin{Verbatim}
PROCEDURE Swap(BYREF a : INTEGER, BYREF b : INTEGER)
    DECLARE temp : INTEGER
    temp ← a
    a ← b
    b ← temp
ENDPROCEDURE
\end{Verbatim}

\subsection*{Local vs global variables}

A \textbf{local variable} 局部变量 is declared inside a subroutine and exists only while it runs. A \textbf{global variable} 全局变量 is declared outside and is visible everywhere. Prefer locals and parameters — heavy use of globals makes code hard to follow and test. (The region where a name is visible is its \textbf{scope} 作用域.)

\subsection*{When to use a subroutine}

Use one when the same logic appears in \textbf{more than one place}, when a block has a \textbf{clear named purpose}, when the program is \textbf{complex} (use \textbf{decomposition} 分解), or when you want to \textbf{test a piece in isolation}. Don't make them so tiny that the call costs more than the content.

\subsection*{Terminology}

\begin{itemize}
\item \textbf{definition} — the \texttt{PROCEDURE ... ENDPROCEDURE} (or function) block.

\item \textbf{call} — where it is invoked. \textbf{argument} — a value passed in. \textbf{parameter} — the variable that receives it.

\item \textbf{return value} — what a function passes back.

\item \textbf{signature} 签名 — name + parameters + return type.

\end{itemize}

\section*{Writing efficient pseudocode}

\begin{itemize}
\item \textbf{move invariants out of loops} — if a value (an \textbf{invariant} 不变量) does not change with the loop counter, compute it once before the loop.

\item \textbf{exit a loop early} when the answer is found (stop a \textbf{linear search} 线性查找 as soon as the target appears).

\item \textbf{avoid redundant work} — store a result and reuse it instead of recomputing.

\item \textbf{choose the right data structure} — an array beats many separate variables when the items belong together.

\item \textbf{replace deep nested IFs with CASE} when testing one value against many.

\item \textbf{comment the intent}, not the mechanics (\texttt{// validate the postcode}, not \texttt{// loop 6 times}).

\item \textbf{use meaningful names} (\texttt{numberOfPupils}, not \texttt{n}) and \textbf{initialise variables} before use.

\end{itemize}


\end{document}
